% =====================================================================
%  THÈME 1 — Pondération des équations chimiques
%  Niveau DIFFICILE — Exercices
% =====================================================================
\documentclass[11pt,a4paper]{article}
\usepackage[utf8]{inputenc}
\usepackage[T1]{fontenc}
\usepackage{lmodern}
\usepackage[french]{babel}
\usepackage{amsmath,amssymb,mathtools}
\usepackage{icomma}
\usepackage[version=4]{mhchem}
\usepackage{siunitx}
\sisetup{output-decimal-marker={,},per-mode=symbol}
\usepackage{xcolor}
\usepackage{tikz}
\usepackage[most]{tcolorbox}
\usepackage{enumitem}
\usepackage{array}
\usepackage{booktabs}
\usepackage{fancyhdr}
\usepackage[a4paper,margin=2.1cm,headheight=26pt,headsep=8pt]{geometry}

\definecolor{primary}{RGB}{142,93,168}
\definecolor{primarydark}{RGB}{82,45,108}
\definecolor{difcol}{RGB}{176,40,60}

\pagestyle{fancy}\fancyhf{}
\fancyhead[L]{\small\textcolor{primarydark}{\textbf{Fiche 5} — Thème 1 : Pondération}}
\fancyhead[R]{\small\textcolor{difcol}{\textsc{Difficile}}}
\fancyfoot[C]{\small\thepage}
\renewcommand{\headrulewidth}{0.8pt}
\renewcommand{\headrule}{\hbox to\headwidth{\color{primary}\leaders\hrule height \headrulewidth\hfill}}

\newtcolorbox[auto counter]{exo}[1][]{enhanced,breakable,boxrule=1pt,arc=3pt,
  left=3mm,right=3mm,top=2.4mm,bottom=2.4mm,colback=difcol!5,colframe=difcol,
  fonttitle=\bfseries,coltitle=white,
  title={Exercice~\thetcbcounter\ \ifx\relax#1\relax\relax\else\textmd{--- \itshape #1}\fi}}

\setlength{\parindent}{0pt}
\setlength{\parskip}{3pt}

\begin{document}

\begin{center}
{\LARGE\bfseries\color{primarydark}Thème 1 — Pondération}\\[1pt]
{\large\color{difcol}\textbf{Niveau Difficile}}\\[2pt]
{\itshape Objectif : pondérer des équations complexes (longues chaînes, ions, redox) par étapes.}\\
{\color{primary}\rule{0.6\linewidth}{1pt}}
\end{center}

\begin{exo}[combustion du kérosène \ce{C14H30}]
Le kérosène (carburant d'avion) est modélisé par \ce{C14H30}. Sa combustion complète s'écrit :
\[ \ce{C14H30 + O2 -> CO2 + H2O}\qquad\text{(non pondérée).} \]
Construis la pondération \textbf{étape par étape} :
\begin{enumerate}[label=\textbf{\alph*)},leftmargin=2em,topsep=1pt,itemsep=1.5pt]
  \item Combien de \ce{CO2} le carbone impose-t-il à droite ?
  \item Combien de \ce{H2O} l'hydrogène impose-t-il à droite ?
  \item Compte alors les atomes d'oxygène à droite. Quel coefficient \emph{fractionnaire} obtiens-tu
        devant \ce{O2} ?
  \item Reviens à des coefficients \textbf{entiers} et écris l'équation finale.
  \item Vérifie C, H et O, puis donne la \textbf{somme} des coefficients.
\end{enumerate}
\end{exo}

\begin{exo}[vérifier et corriger une équation ionique]
On propose l'équation de la réaction du cobalt avec le nitrate de potassium en milieu chlorhydrique :
\[ \ce{3 Co + 2 KNO3 + 8 HCl -> 2 NO + 3 CoCl2 + 1 KCl + 4 H2O}\quad(?) \]
\begin{enumerate}[label=\textbf{\alph*)},leftmargin=2em,topsep=1pt,itemsep=1.5pt]
  \item Dresse l'inventaire complet des atomes (\ce{Co}, \ce{K}, \ce{N}, \ce{O}, \ce{H}, \ce{Cl}),
        réactifs \emph{vs} produits.
  \item Quels éléments ne sont \textbf{pas} équilibrés ?
  \item Corrige l'équation en changeant \emph{un seul} coefficient. Écris l'équation ajustée.
  \item Donne la somme des coefficients de l'équation correcte.
\end{enumerate}
\end{exo}

\begin{exo}[une équation redox « moléculaire » à plusieurs groupements]
La réduction du dichromate de potassium par le carbone en milieu sulfurique s'écrit :
\[ \ce{K2Cr2O7 + C + H2SO4 -> CO2 + Cr2(SO4)3 + K2SO4 + H2O}\qquad\text{(non pondérée).} \]
Pondère par étapes :
\begin{enumerate}[label=\textbf{\alph*)},leftmargin=2em,topsep=1pt,itemsep=1.5pt]
  \item Équilibre d'abord le \textbf{chrome} (Cr) et le \textbf{potassium} (K) en fixant les
        coefficients de \ce{K2Cr2O7}, \ce{Cr2(SO4)3} et \ce{K2SO4}. \emph{(Astuce : essaie
        $2\,\ce{K2Cr2O7}$.)}
  \item Déduis-en le nombre total de groupements \textbf{\ce{SO4}}, donc le coefficient de \ce{H2SO4}.
  \item Équilibre le \textbf{carbone} (C $\to$ \ce{CO2}), puis l'\textbf{hydrogène} par l'eau.
  \item Écris l'équation finale et vérifie l'\textbf{oxygène} pour contrôle. Donne la somme des coefficients.
\end{enumerate}
\end{exo}

\begin{exo}[équation ionique redox avec charges : dichromate + fer(II)]
En solution acide, les ions dichromate oxydent les ions fer(II) :
\[ \ce{Cr2O7^2- + H3O+ + Fe^2+ -> Cr^3+ + H2O + Fe^3+}\qquad\text{(non pondérée).} \]
\begin{enumerate}[label=\textbf{\alph*)},leftmargin=2em,topsep=1pt,itemsep=1.5pt]
  \item Équilibre le \textbf{chrome} (fixe le coefficient de \ce{Cr^3+}).
  \item On admet les coefficients $14$ devant \ce{H3O+}, $6$ devant \ce{Fe^2+}, $21$ devant \ce{H2O}
        et $6$ devant \ce{Fe^3+}. Vérifie que l'\textbf{oxygène} et l'\textbf{hydrogène} sont équilibrés.
  \item Vérifie que la \textbf{charge totale} est identique des deux côtés (c'est la condition décisive
        d'une équation ionique). Conclus.
\end{enumerate}
\end{exo}

\end{document}
